%=============================================================================
%  11_reporting.tex  -  Chapter 10: Cleaning and PDF Reporting
%=============================================================================
\chapter{Cleaning and PDF Reporting}
\label{ch:reporting}

Two skills turn analysis into deliverables: the cleaning skill, which produces a
tidier dataset with a transparent record of what changed, and the report skill,
which exports a branded PDF capturing the whole session. Both are held to the
same honesty standard as the rest of the system.

\section{The cleaning skill}
The cleaning skill applies guided data-cleaning operations to a session's
DataFrame and returns a structured \code{CleaningResponse}. Critically, it does
not clean silently: the response enumerates exactly what happened.

\begin{table}[h]
\centering
\caption{Fields of \code{CleaningResponse}.}
\small
\begin{tabularx}{\linewidth}{@{}L{3.6cm} X@{}}
\toprule
\textbf{Field} & \textbf{Meaning} \\
\midrule
\code{original\_rows} / \code{cleaned\_rows} & Row counts before and after. \\
\code{original\_columns} / \code{cleaned\_columns} & Column counts before and
after. \\
\code{actions} & An ordered list of the cleaning operations performed. \\
\code{explanation} & A human-readable rationale for the changes. \\
\code{profile} & A fresh \code{DataProfile} of the cleaned data. \\
\code{preview\_rows} & A sample of the cleaned rows for immediate inspection. \\
\bottomrule
\end{tabularx}
\end{table}

Because the response includes a fresh profile and a preview, the front end can
show the effect of cleaning immediately, and the user can decide whether to keep
the result --- optionally as a new session, via the
\code{create\_new\_session} flag on the request. Transparency here is the point:
an analyst must be able to see and justify every transformation applied to their
data.

\begin{uanote}[Cleaning is auditable]
Every cleaning run returns the list of actions taken and a before/after profile.
There is no ``magic tidy'' that changes the data without saying how --- the
transformation is as inspectable as the analysis.
\end{uanote}

\section{The report skill}
The report skill produces a professional, branded PDF using \textbf{ReportLab}.
It is deliberately designed so that ``the PDF and the app read as one product'':
the same indigo/amber identity, the same monospaced treatment of numbers, and
the same restraint.

\subsection{Document structure}
A report is assembled with ReportLab's \code{platypus} flowable system on A4
pages. A custom stylesheet defines a coherent hierarchy --- title, subtitle,
eyebrow, headings, body, muted, and small styles --- so the document has
deliberate typographic rhythm rather than default styling.

\begin{table}[h]
\centering
\caption{Sections composed into the PDF report.}
\small
\begin{tabularx}{\linewidth}{@{}L{3.4cm} X@{}}
\toprule
\textbf{Section} & \textbf{Content} \\
\midrule
Cover / header & Branded title block establishing the report identity. \\
Data profile & Row/column counts, per-column types, and key statistics. \\
Quality signals & Missingness, outliers, and other profile-derived signals. \\
Generated charts & The validated charts produced during the session. \\
Narrative & A model-written summary from the report-narrative generator. \\
Transcript & The conversation that produced the analysis. \\
\bottomrule
\end{tabularx}
\end{table}

\subsection{Embedding real charts}
The charts in the report are the very same validated figures generated during
the session --- not re-imagined for print. Plotly figures are rasterized to
images through \textbf{Kaleido} and embedded into the PDF, so what the analyst
saw on screen is what appears in the document. This closes the loop: a chart
that passed validation in the dashboard is the chart that lands in the report.

\subsection{The report narrative}
A dedicated \code{report\_narrative} generator asks the model for a concise,
prose summary of the findings, which is embedded as the narrative section. It
uses the non-streaming \code{generate\_text} path (Chapter~\ref{ch:llm}), so a
model failure raises rather than silently producing an empty section --- once
again, no fabrication.

\begin{lstlisting}[style=py,caption={ReportLab flowables used to build the document.}]
from reportlab.lib.pagesizes import A4
from reportlab.lib.styles import ParagraphStyle, getSampleStyleSheet
from reportlab.platypus import (SimpleDocTemplate, Paragraph, Spacer,
                                Table, TableStyle, Image)
# A custom stylesheet (title/subtitle/eyebrow/h1/h2/body/muted/small)
# gives the PDF the same visual identity as the application.
\end{lstlisting}

\subsection{Robust value formatting}
Report generation is defensive about the messy reality of real data. Helper
routines (\code{\_fmt}, \code{\_plain}, \code{\_safe\_text},
\code{\_inline\_markdown}, \code{\_top\_value\_text}) coerce arbitrary profile
values into clean, length-bounded strings, strip inline Markdown, and format
top-value lists --- so a stray \code{NaN}, an over-long category label, or an
exotic type never breaks the layout.

\section{Delivery}
The report is requested through \code{POST /skills/report}. In the chat flow, a
report intent additionally emits a \code{download\_report} command event so the
front end can trigger the download as part of the conversation, making export
feel like a natural continuation of the analysis rather than a separate mode.
